We now develop the formal model: the change universe, the acceptance predicate,
the state transition law, and the membership theorem, followed by the threat
model under which the results of Sections~\ref{sec:inaccessibility}
and~\ref{sec:attacks} hold.

\subsection{The Change Universe and Admitted Set}

Let $\mathcal{A}$ be a set of actors, $\mathcal{E}$ a set of event types,
$\mathcal{O}$ a set of typed object-relation sets (each element an OCEL~2.0
event-to-object relation), $\mathcal{S}$ a set of process states, $\mathcal{P}$
a set of policy versions, $\mathcal{V}$ a set of validator versions,
$\mathcal{T}$ a set of timestamps, $\mathcal{N}$ a set of nonces, $\mathcal{R}$
a set of receipts, and $\Sigma$ a set of signatures.

\begin{definition}[Enterprise Change Universe]
\label{def:omega}
The \emph{enterprise change universe} is
\[
  \Omega_{\mathrm{change}} \;=\; \mathcal{A} \times \mathcal{E} \times
  \mathcal{O} \times \mathcal{S} \times \mathcal{R} \times \mathcal{P} \times
  \mathcal{V} \times \mathcal{T} \times \mathcal{N} \times \mathcal{R} \times
  \Sigma.
\]
A \emph{candidate change} is a tuple
$x = (a, e, O, s, R_{\mathrm{prior}}, p, v, t, n, R, \sigma)$: actor $a$
proposes event $e$ over object relations $O$ from process state $s$, citing
prior receipt $R_{\mathrm{prior}}$, under policy version $p$ and validator
version $v$, at time $t$ with freshness nonce $n$, carrying output receipt $R$
and signature $\sigma$.
\end{definition}

The universe is deliberately maximal: \emph{every} mutation an enterprise can
undergo --- data edits, status updates, approvals, payments, permission
changes, deployments, configuration and policy updates, document mutations,
identity changes, model updates, and AI-agent actions --- is representable as a
candidate change. The framework's discipline consists in what it
\emph{accepts} from this universe, not in what can be expressed.

In the running P2P example, the legitimate approval of invoice \#4711 is the
candidate
\begin{align*}
  x_{\textsf{appr}} = (&\textsf{emp\_207},\, \textsf{Approve},\,
  \{\textsf{Inv\#4711}, \textsf{PO\#982}\},\, \textsf{submitted},\\
  &R_{17},\, p_{3},\, v_{9},\, t,\, n,\, R_{18},\, \sigma),
\end{align*}
where $R_{17}$ is the receipt of the preceding \textsf{Submit} and $R_{18}$
the receipt this approval emits.

\begin{definition}[Valid Receipt]
\label{def:valid-receipt}
A receipt $R$ is \emph{valid} for the transition $s \xrightarrow{e} s'$,
written $\verifyreceipt(R) = 1$, iff all of the following hold:
\begin{enumerate}[label=(\roman*),leftmargin=2.4em,noitemsep]
  \item \emph{Substance:} $R$ embeds, or references by content address, the
        complete expected and observed OCEL~2.0 logs of the transition, in
        canonical serialisation, such that any holder can recompute every hash
        in $R$;
  \item \emph{Chaining:} $R$ carries $h(R_{\mathrm{prior}})$,
        $h(\textit{in})$, and $h(\textit{out})$, and $h(R_{\mathrm{prior}})$
        equals the hash of the current chain head;
  \item \emph{Verifier authority:} the alignment state in $R$ was computed and
        sealed by the verifier; producer-asserted alignment is rejected;
  \item \emph{Challenge binding:} a nonce minted before execution appears in
        the challenge manifest, in the raw runner evidence, and inside the
        observed OCEL events themselves.
\end{enumerate}
\end{definition}

Condition (i) is the formal content of the axiom of receipt substance
(Section~\ref{sec:bg-receipts}); conditions (iii) and (iv) close the two
historically exploited shortcuts --- self-certification and fixture replay ---
that Section~\ref{sec:impl-refusals} enumerates as refusal states.

\begin{definition}[Acceptance Predicate]
\label{def:accept}
For a candidate change
$x = (a, e, O, s, R_{\mathrm{prior}}, p, v, t, n, R, \sigma)$ as in
Definition~\ref{def:omega}, define
\begin{align}
  \accept(x) \;=\;\;
  &\verifysig\bigl(\sigma,\, h(a \| e \| O \| s \| R_{\mathrm{prior}} \| p \| v \| t \| n)\bigr) \tag{C1}\\
  \wedge\;\; &\verifyreceipt(R) \tag{C2}\\
  \wedge\;\; &\operatorname{PolicyAllowed}(a, e, p) \tag{C3}\\
  \wedge\;\; &\operatorname{ValidatorAllowed}(v) \tag{C4}\\
  \wedge\;\; &\operatorname{Fresh}(n) \tag{C5}\\
  \wedge\;\; &B(s, e) = 1 \tag{C6}\\
  \wedge\;\; &\operatorname{ObjectsValid}(O). \tag{C7}
\end{align}
$\operatorname{Fresh}(n)$ holds iff $n$ has not been consumed by any previously
accepted change. $\operatorname{PolicyAllowed}(a,e,p)$ holds iff policy version
$p$ is current and grants actor $a$ the authority to emit event type $e$.
$\operatorname{ValidatorAllowed}(v)$ holds iff $v$ identifies a non-revoked
verifier build. $\operatorname{ObjectsValid}(O)$ holds iff the event-to-object
relation $O$ is well-typed against the OCEL~2.0 schema and consistent with the
object lifecycles recorded in the receipt chain.
\end{definition}

Each conjunct is independently checkable and independently falsifiable; this
independence is what powers the decomposition of Section~\ref{sec:attacks}.
Note also the layering: C1 and C2 are cryptographic, C3--C5 are operational
bindings, C6 is the process law, and C7 is the object law.

\begin{definition}[Admitted Change Set]
\label{def:A}
$\Aent \;=\; \{\, x \in \Omega_{\mathrm{change}} \mid \accept(x) = 1 \,\}$.
\end{definition}

\begin{theorem}[Acceptance Implies Membership]
\label{thm:main}
For every $x \in \Omega_{\mathrm{change}}$:
$\accept(x) = 1 \Rightarrow x \in \Aent$.
\end{theorem}
\begin{proof}
Immediate from Definition~\ref{def:A}: $\Aent$ is the preimage
$\accept^{-1}(1)$.
\end{proof}

Theorem~\ref{thm:main} is definitionally true, and we state it anyway, because
its force is architectural rather than mathematical: it asserts that
\emph{there exists no second path} into enterprise state. In conventional
architectures the analogous statement is false --- state is reachable through
the ORM, through the admin panel, through the migration script, through the
integration webhook --- and every such path is an acceptance gate that the
process model never sees. The theorem is the design obligation that all of
those paths either emit admitted changes or lose state authority.

\begin{corollary}[Refusal or Residualization]
\label{cor:refuse}
For every $x \in \Omega_{\mathrm{change}}$:
$x \notin \Aent \Rightarrow \refuse(x) \vee \residualize(x)$.
\end{corollary}

$\refuse(x)$ rejects the change outright, returning the failing conjunct.
$\residualize(x)$ quarantines it: the candidate is recorded, content-addressed,
and preserved as investigative evidence, but acquires no state authority.
Residualization is the lossless refusal mode --- it matters operationally
because enterprises cannot simply drop nonconforming inputs (a malformed
payment instruction is itself evidence of something), and it matters for the
process-mining feedback loop: clusters of similar residuals indicate either an
attack campaign or a legitimate behaviour the declared model is missing, in
which case the remedy is an explicit, reviewed extension of the process law ---
never a silent widening of the gate.

\begin{definition}[State Transition Law]
\label{def:mu}
The enterprise state transition function $\mu : \mathcal{S} \times \mathcal{E}
\to \mathcal{S}$ satisfies
\[
  \mu(s, e) \;=\;
  \begin{cases}
    s' & \text{if } B(s,e) = 1 \text{ and } \accept(x) = 1
         \text{ for the change } x \text{ carrying } e,\\[2pt]
    s  & \text{otherwise.}
  \end{cases}
\]
\end{definition}

\begin{lemma}[Chain Soundness]
\label{lem:chain}
If every accepted change satisfies Definition~\ref{def:accept}, then for any
reachable state $s_n$ there exists a receipt chain
$R_0 \to \cdots \to R_{n-1}$ such that
$s_n = \mu(\mu(\cdots\mu(s_0, e_0)\cdots), e_{n-1})$ and each link is
independently re-verifiable from the receipt substance alone.
\end{lemma}
\begin{proof}
By induction on $n$. The base case is the genesis receipt $R_0$ witnessing
$s_0$. For the step, the change accepted at position $t$ satisfies C2, so
$R_t$ embeds the canonical logs and chains to $h(R_{t-1})$
(Definition~\ref{def:valid-receipt}, i--ii); by the induction hypothesis the
prefix is re-verifiable, and condition (i) makes the new link re-verifiable
without reference to any state outside the chain.
\end{proof}

Lemma~\ref{lem:chain} is what converts the framework from an access-control
discipline into an \emph{evidence} discipline: current state is not a row that
happens to exist but the endpoint of a replayable derivation. It is also the
precondition for retrospective process mining over admitted history ---
discovery and conformance algorithms~\cite{van_der_aalst_2016_process_mining}
run over the chain's embedded OCEL logs with the unusual guarantee that the log
is complete and tamper-evident by construction.

\subsection{Threat Model}
\label{sec:threat-model}

We state the adversary explicitly. The adversary \textbf{may}: submit arbitrary
candidate changes at any rate; read all public state, receipts, and policies;
replay any previously observed message; possess valid credentials for some
actors (insider); hold administrative privilege in surrounding infrastructure
(database, message bus, file system); and operate AI agents with write-capable
tools. The adversary \textbf{may not} (assumption ledger, cf.\
Section~\ref{sec:attacks}): forge signatures under EUF-CMA of the deployed
scheme; find BLAKE3 preimages or collisions; compromise the validator's signing
key; or exploit a defect in the gate implementation. Each ``may not'' is a
named channel in Theorem~\ref{thm:decomp}, with its own probability term and
mitigation --- the framework's security claim is exactly the conjunction of
these assumptions, no more.

Two consequences of the threat model deserve emphasis. First, database
administrators are \emph{inside} the adversary model: a raw write by a DBA is a
candidate change like any other and fails C2
(Proposition~\ref{prop:raw}). The framework therefore does not trust the
storage layer --- corrupted rows can exist, but they are detectable as
divergence between stored state and the receipt-chain derivation of
Lemma~\ref{lem:chain}. Second, the verifier is trusted but \emph{versioned and
revocable} (C4): a compromised validator build is contained by revoking $v$,
which retroactively identifies every change it admitted.

\subsection{The Chesterton Fence as a State Law}

Chesterton's principle~\cite{chesterton_orthodoxy_1908} holds that a fence
should not be removed until its purpose is understood. The boundary predicate
$B(s,e)$ is this principle expressed as a transition law: every enterprise
transition is fenced, the fence is versioned and signed, and crossing it
requires the full witness of Definition~\ref{def:accept}. Crucially, the fence
does not prevent the \emph{attempt} --- the adversary may submit any
$x \in \Omega_{\mathrm{change}}$ --- it prevents the attempt from
\emph{counting}. And symmetrically, the fence cannot be removed silently:
relaxing $B$ is itself a policy change, hence itself an admitted change with
its own receipt, reviewable in the chain like any other.
